\section{A quantitative prime-length bond-dimension dichotomy}

Fix one positive local interaction. A sufficiently small matrix product
representation at one prime length determines a single tensor generating all
sufficiently long ground spaces. Otherwise every prime-length representation
obeys an explicit fifth-degree obstruction. This is the quantified asymptotic
form of \cite[Theorem~5]{Sanz2010Quantum}; the finite-prefix boundary is
recorded in \cite{gap:spwc10_ground_state_dichotomy_quantifiers}.

\begin{lemma}[Normal blocks at lengths without small divisors]
    \label{lem:normal_blocks_no_small_divisor}
    \lean{MPSTensor.exists_normal_blocks_of_no_small_divisor}
    \leanok
    \uses{def:irreducible_tensor, def:l_blk_injective,
        def:blocks_not_gauge_phase_equiv}
    Let $A$ have physical dimension $d$, bond dimension $D$, and chain length
    $N\ge2$. Suppose no integer $m$ with $2\le m\le D$ divides $N$. There
    exist $b\le D$ blocks $B_j$ of positive bond dimensions $D_j\le D$ and
    coefficients $c_j\ne0$ such that
    \begin{align}
        V^{(N)}(A) & =\sum_{j=1}^b c_jV^{(N)}(B_j).
        \label{eq:normal_blocks_no_small_divisor}
    \end{align}
    Every $B_j$ is unital, is irreducible, has a positive definite dual fixed
    point, and is injective at the common length $q(D)=(D^2+1)^2$.
    Distinct retained blocks are inequivalent under a gauge and a phase.
    The decomposition and all these properties are derived from $A$.
\end{lemma}
\begin{proof}\leanok
    \uses{thm:pgvwc07_positive_length_witness, lem:pgvwc07_block_irreducible,
        lem:periodic_scale_of_irreducible, thm:periodic_period_le_bondDim,
        thm:pgvwc07_normalized_periodic_state_vector_decomposition,
        thm:wielandt_general, lem:gauge_invariance_block_injective,
        thm:normality_rescaling_invariance, def:mpv_phase_class_data,
        thm:mpv_decomposition}
    Lemma~\ref{thm:pgvwc07_positive_length_witness} reduces $A$ to unital
    canonical blocks $A^k$ with positive definite dual fixed points, weights
    $\mu_k$, and $\sum_k\dim A^k\le D$; in particular there are at most $D$
    blocks. By Lemma~\ref{lem:pgvwc07_block_irreducible} every block is
    irreducible, and Lemma~\ref{lem:periodic_scale_of_irreducible} rescales
    and gauges it to a periodic tensor of period $m_k$, whose period is at
    most $D$ by Lemma~\ref{thm:periodic_period_le_bondDim}. If $m_k\ne1$,
    then $2\le m_k\le D$, so $m_k\nmid N$ and
    Theorem~\ref{thm:pgvwc07_normalized_periodic_state_vector_decomposition}
    gives $V^{(N)}(A^k)=0$. If $m_k=1$, the representative is primitive, and
    the general Wielandt bound of Theorem~\ref{thm:wielandt_general} gives
    injectivity at $q(D)$; it transfers to $A^k$ through
    Lemma~\ref{lem:gauge_invariance_block_injective} and
    Lemma~\ref{thm:normality_rescaling_invariance}. Group the blocks into
    the phase classes of Definition~\ref{def:mpv_phase_class_data}. The
    $s$-th member $A^{(j,s)}$ of class $j$, with weight $\mu_{j,s}$, satisfies
    $V^{(N)}(A^{(j,s)})=\xi_{j,s}^NV^{(N)}(A^{(j,0)})$ for some
    $\xi_{j,s}\ne0$. Put $\gamma_j=\sum_s\mu_{j,s}^N\xi_{j,s}^N$. By
    Theorem~\ref{thm:mpv_decomposition},
    \begin{align}
        V^{(N)}(A) & =\sum_k\mu_k^NV^{(N)}(A^k)
        =\sum_j\gamma_jV^{(N)}(A^{(j,0)}).
    \end{align}
    Discard the classes with $\gamma_j=0$ and those whose representative
    has period other than one, whose vectors vanish. The remaining
    representatives and coefficients are the blocks $B_j$ and the
    coefficients $c_j$ of~(\ref{eq:normal_blocks_no_small_divisor}); no two
    of them are related by a gauge transformation and a phase.
\end{proof}

\begin{lemma}[One small prime-length representation gives a fixed eventual tensor]
    \label{lem:prime_small_representation_fixed_tensor}
    \lean{MPSTensor.exists_normal_eventual_ground_data_of_prime_length_bound,
        MPSTensor.exists_eventually_fixed_mps_of_prime_length_bound}
    \leanok
    \uses{lem:normal_blocks_no_small_divisor,
        lem:finite_volume_local_kernel_inclusion, lem:local_mps_kernel_all_volumes}
    Let $h\ge0$ act on $L$ sites, and assume $\dim\ker H_M=1$ for every
    $M\ge L$. Suppose $N$ is prime, $V^{(N)}(A)\ne0$ is annihilated by
    $H_N$, and
    \begin{align}
        N & \ge L+\max\{q(D),3(D-1)(q(D)+1)\},
          & q(D)                               & =(D^2+1)^2.
        \label{eq:prime_small_representation_threshold}
    \end{align}
    There exist $0<D_0\le D$, one tensor $B$ of bond dimension $D_0$, and
    $M_0$ such that, for every $M\ge M_0$,
    \begin{align}
        V^{(M)}(B) & \ne0,
                   & \ker H_M & =\C V^{(M)}(B).
        \label{eq:prime_small_representation_fixed_tensor}
    \end{align}
\end{lemma}
\begin{proof}\leanok
    \uses{lem:normal_blocks_no_small_divisor,
        lem:finite_volume_local_kernel_inclusion, lem:local_mps_kernel_all_volumes}
    The threshold gives $D<N$, so primality permits
    Lemma~\ref{lem:normal_blocks_no_small_divisor}. The original nonzero vector
    ensures a retained block exists. Simultaneous complementary spanning
    and Lemma~\ref{lem:finite_volume_local_kernel_inclusion} give
    $G_L(B_j)\subseteq\ker h$ for each retained block. Choose any one.
    Its periodic vector has zero energy at all permitted volumes and is
    nonzero for $M\ge q(D)+1$. The one-dimensional kernel then equals its line.
    One may take $M_0=\max\{L,q(D)+1\}$.
\end{proof}

\begin{theorem}[Fixed eventual tensor or a prime-length fifth-degree obstruction]
    \label{thm:eventual_mps_prime_length_dichotomy}
    \lean{MPSTensor.eventually_fixed_mps_or_prime_length_bound}
    \leanok
    \discussion{10280}
    \uses{lem:prime_small_representation_fixed_tensor}
    Let $h\ge0$ be one fixed $L$-site interaction, and assume
    $\dim\ker H_M=1$ for every $M\ge L$. Then either one fixed tensor $B$
    generates the nonzero ground-state line at every sufficiently large
    volume, or every prime $N$ and every bond-$D$ tensor $A$ with
    $0\ne V^{(N)}(A)\in\ker H_N$ satisfy
    \begin{align}
        N<L+\max\{(D^2+1)^2,3(D-1)((D^2+1)^2+1)\}.
        \label{eq:eventual_mps_prime_length_dichotomy}
    \end{align}
    The interaction and $L$ are fixed independently of $N,D,A$.
    \textbf{Scope restriction (eventual uniform branch):}
    the theorem supplies the same tensor at all sufficiently large volumes;
    the finite-prefix representation argument for a literal all-volume
    formulation is separate; see
    \cite{gap:spwc10_ground_state_dichotomy_quantifiers}.
\end{theorem}
\begin{proof}\leanok
    \uses{lem:prime_small_representation_fixed_tensor}
    If the eventual fixed-tensor alternative fails, the contrapositive of
    Lemma~\ref{lem:prime_small_representation_fixed_tensor} gives the displayed
    strict inequality for every prime-length nonzero ground-state representation.
\end{proof}

\subsection{Finite prefixes and exact cyclic momentum}
\label{sec:ground_state_finite_prefix}

\begin{lemma}[Periodic representation of an exactly invariant finite vector]
    \label{lem:periodic_finite_cyclic_representation}
    \lean{MPSTensor.exists_mpv_eq_of_cyclic_invariant,
        MPSTensor.exists_mpvState_eq_of_cyclic_invariant,
        MPSTensor.mpvState_cyclic_invariant,
        MPSTensor.exists_mpvState_eq_iff_cyclic_invariant}
    \notready
    \uses{thm:seqanc_cut_rank_attained, thm:seqanc_cut_rank_bounds,
        thm:pgvwc07_site_independent_matrices}
    Let $N>0$ and let $\psi\ne0$ be a finite vector satisfying
    $\psi(\sigma\circ\tau_s)=\psi(\sigma)$ for every cyclic translation
    $\tau_s$. There is a periodic trace-MPS representation of $\psi$
    with matrix size $N d^{\lfloor N/2\rfloor}$. Conversely every
    periodic trace-MPS vector has this exact cyclic invariance.
    Thus a nonzero positive-length vector is representable precisely
    when its cyclic momentum is zero.
\end{lemma}
\begin{proof}
    The actual finite OBC construction supplies a chain with bond bound
    $d^{\lfloor N/2\rfloor}$. Apply the cyclic block averaging of
    \cite[Theorem~3]{PerezGarcia2007Matrix}; exact invariance makes the
    average equal to the original vector. The converse is cyclicity
    of the matrix trace, including every power of the basic cyclic shift.
\end{proof}

\begin{lemma}[Unused periodic bond coordinates]
    \label{lem:periodic_uniform_bond_padding}
    \lean{MPSTensor.padBond, MPSTensor.mpv_padBond, MPSTensor.mpvState_padBond}
    \notready
    \uses{def:mps_tensor}
    Extending all site matrices by zero from dimension $D$ to any
    dimension $K\ge D$ preserves every periodic vector of positive
    length. The assertion excludes the empty word, whose trace is the
    bond dimension.
\end{lemma}
\begin{proof}
    Zero extension preserves products through a retained intermediate
    bond. Induct on nonempty words and take the trace; all added
    diagonal entries are zero.
\end{proof}

\begin{lemma}[One common bond size on the exact cyclic ground sector]
    \label{lem:uniform_bond_cyclic_ground_sector}
    \lean{MPSTensor.exists_uniform_bond_on_cyclic_ground_sector}
    \notready
    \uses{thm:eventual_mps_prime_length_dichotomy,
        lem:periodic_finite_cyclic_representation, lem:periodic_uniform_bond_padding}
    Suppose the permitted ground spaces $\ker H_M$, $M\ge L$, are
    one-dimensional, and a bond-$D$ tensor $B$ generates their nonzero
    ground lines for every $M\ge M_0$. Then the single matrix size
    \[
        K=D+M_0d^{M_0}+1
    \]
    suffices to generate each permitted positive-length ground line
    containing a nonzero exactly cyclic-invariant vector, including
    every such line at $M<M_0$. The representing tensor may depend
    on $M$.
\end{lemma}
\begin{proof}
    At long lengths extend $B$ by unused bond coordinates. At a short
    length use the preceding finite representation, whose dimension is
    $M d^{\lfloor M/2\rfloor}\le M_0d^{M_0}<K$, and then extend it.
    The actual nonzero ground vector and uniqueness identify the whole
    ground line.
\end{proof}

\begin{lemma}[Exact finite-prefix criterion for an all-volume family]
    \label{lem:uniform_bond_finite_prefix_criterion}
    \lean{MPSTensor.uniform_bond_family_iff_short_cyclic_ground}
    \notready
    \uses{lem:uniform_bond_cyclic_ground_sector}
    Under the preceding assumptions, the following are equivalent:
    \begin{enumerate}
        \item There is one positive matrix size $K$ such that, at
              every permitted positive volume $M$, some bond-$K$ tensor
              has a nonzero periodic vector generating $\ker H_M$.
        \item At each permitted short positive volume $M<M_0$, the
              ground space contains a nonzero vector fixed exactly by
              every cyclic translation.
    \end{enumerate}
    An invariant ground line alone does not state the second condition:
    translation can act on a line by a nontrivial root of unity.
    No derivation of the second condition from the complete hypotheses
    of \cite[Theorem~5]{Sanz2010Quantum} is asserted here; this remaining
    boundary is recorded in
    \cite{gap:spwc10_ground_state_dichotomy_quantifiers}.
\end{lemma}
\begin{proof}
    The first condition supplies nonzero periodic trace vectors, and
    trace cyclicity proves the second. Conversely use the preceding
    common-size construction at the short volumes and the already
    supplied eventual tensor at the long volumes.
\end{proof}

